\begin{figure}[tb]
\centering
\begin{tikzpicture}[scale=1.05,>=Stealth]
 \draw[gray!60,thin] (1.55,0) circle[radius=1.55];
 \draw[black,thick] (2,0) circle[radius=2];
 \draw[gray!75,dashed] (2.45,0) circle[radius=2.45];
 \draw[->,gray] (-.35,0)--(5.35,0) node[right] {\(\theta\)};
 \fill (0,0) circle (1.8pt) node[below left] {\(z\)};
 \fill (1.55,0) circle (1.3pt) node[below=5pt] {\(z+(l-v)\theta\)};
 \fill (2,0) circle (1.6pt) node[above=4pt] {\(x=z+l\theta\)};
 \fill (2.45,0) circle (1.3pt);
 \draw[gray] (2.45,-.07)--(5.3,-.85)
   node[right,black] {\(z+(l+v)\theta\)};
 \fill (2,2) circle (2pt) node[above right] {\(y\)};
 \draw[dotted] (0,0)--(2,2);
 \node[align=left,anchor=west] at (5.3,1.65)
   {\(\displaystyle B(z+L\theta,L)\)\\[2pt]increases with \(L\)};
\end{tikzpicture}
\caption{The tangent-disk sweep. All boundary circles pass through \(z\).
A shell point \(y\) with \(|y-z|\ge\rho\) enters between the parameters
\(l-v\) and \(l+v\), with \(v\) proportional to
\(\delta/\rho^2\). Mass outside \(B(z,\rho)\) already present at
\(l-v\) lies in the inner disk \(B(x,l-2\delta)\). When
\(z\in D(x,r,\delta)\), this disk is contained in the controlled disk
\(B(x,r-\delta)\). The drawing is schematic.}
\label{fig:tangent-disks}
\end{figure}
